% Options for packages loaded elsewhere
\PassOptionsToPackage{unicode}{hyperref}
\PassOptionsToPackage{hyphens}{url}
\documentclass[
  ignorenonframetext,
]{beamer}
\newif\ifbibliography
\usepackage{pgfpages}
\setbeamertemplate{caption}[numbered]
\setbeamertemplate{caption label separator}{: }
\setbeamercolor{caption name}{fg=normal text.fg}
\beamertemplatenavigationsymbolsempty
% remove section numbering
\setbeamertemplate{part page}{
  \centering
  \begin{beamercolorbox}[sep=16pt,center]{part title}
    \usebeamerfont{part title}\insertpart\par
  \end{beamercolorbox}
}
\setbeamertemplate{section page}{
  \centering
  \begin{beamercolorbox}[sep=12pt,center]{section title}
    \usebeamerfont{section title}\insertsection\par
  \end{beamercolorbox}
}
\setbeamertemplate{subsection page}{
  \centering
  \begin{beamercolorbox}[sep=8pt,center]{subsection title}
    \usebeamerfont{subsection title}\insertsubsection\par
  \end{beamercolorbox}
}
% Prevent slide breaks in the middle of a paragraph
\widowpenalties 1 10000
\raggedbottom
\AtBeginPart{
  \frame{\partpage}
}
\AtBeginSection{
  \ifbibliography
  \else
    \frame{\sectionpage}
  \fi
}
\AtBeginSubsection{
  \frame{\subsectionpage}
}
\usepackage{iftex}
\ifPDFTeX
  \usepackage[T1]{fontenc}
  \usepackage[utf8]{inputenc}
  \usepackage{textcomp} % provide euro and other symbols
\else % if luatex or xetex
  \usepackage{unicode-math} % this also loads fontspec
  \defaultfontfeatures{Scale=MatchLowercase}
  \defaultfontfeatures[\rmfamily]{Ligatures=TeX,Scale=1}
\fi
\usepackage{lmodern}
\ifPDFTeX\else
  % xetex/luatex font selection
\fi
% Use upquote if available, for straight quotes in verbatim environments
\IfFileExists{upquote.sty}{\usepackage{upquote}}{}
\IfFileExists{microtype.sty}{% use microtype if available
  \usepackage[]{microtype}
  \UseMicrotypeSet[protrusion]{basicmath} % disable protrusion for tt fonts
}{}
\makeatletter
\@ifundefined{KOMAClassName}{% if non-KOMA class
  \IfFileExists{parskip.sty}{%
    \usepackage{parskip}
  }{% else
    \setlength{\parindent}{0pt}
    \setlength{\parskip}{6pt plus 2pt minus 1pt}}
}{% if KOMA class
  \KOMAoptions{parskip=half}}
\makeatother
\usepackage{longtable,booktabs,array}
\usepackage{calc} % for calculating minipage widths
\usepackage{caption}
% Make caption package work with longtable
\makeatletter
\def\fnum@table{\tablename~\thetable}
\makeatother
\setlength{\emergencystretch}{3em} % prevent overfull lines
\providecommand{\tightlist}{%
  \setlength{\itemsep}{0pt}\setlength{\parskip}{0pt}}
\usepackage{bookmark}
\IfFileExists{xurl.sty}{\usepackage{xurl}}{} % add URL line breaks if available
\urlstyle{same}
\hypersetup{
  pdftitle={The EM Algorithm for MMLE, Step by Step},
  pdfauthor={Derek Briggs},
  hidelinks,
  pdfcreator={LaTeX via pandoc}}

\title{The EM Algorithm for MMLE, Step by Step}
\subtitle{EDUC 8720: Item Parameter Estimation}
\author{Derek Briggs}
\date{2026-02-15}

\begin{document}
\frame{\titlepage}

\begin{frame}{The EM Algorithm: Overview}
\phantomsection\label{the-em-algorithm-overview}
The EM algorithm alternates between two steps:

\begin{longtable}[]{@{}
  >{\raggedright\arraybackslash}p{(\linewidth - 4\tabcolsep) * \real{0.2143}}
  >{\raggedright\arraybackslash}p{(\linewidth - 4\tabcolsep) * \real{0.4643}}
  >{\raggedright\arraybackslash}p{(\linewidth - 4\tabcolsep) * \real{0.3214}}@{}}
\toprule\noalign{}
\begin{minipage}[b]{\linewidth}\raggedright
Step
\end{minipage} & \begin{minipage}[b]{\linewidth}\raggedright
What it does
\end{minipage} & \begin{minipage}[b]{\linewidth}\raggedright
Analogy
\end{minipage} \\
\midrule\noalign{}
\endhead
\textbf{E-step} & Uses current item parameter guesses to figure out how
examinees are distributed across the ability scale & ``If these were the
true item parameters, where would each examinee likely fall?'' \\
\textbf{M-step} & Uses that ability distribution to re-estimate item
parameters & ``Given where we think examinees fall, what item parameters
best fit the responses?'' \\
\bottomrule\noalign{}
\end{longtable}

We repeat until the estimates stop changing.
\end{frame}

\begin{frame}{The Data: LSAT-6}
\phantomsection\label{the-data-lsat-6}
\begin{itemize}
\tightlist
\item
  \textbf{1000 examinees}, \textbf{5 dichotomous items}, \textbf{32
  unique response patterns}
\end{itemize}

\begin{longtable}[]{@{}lr@{}}
\caption{8 Most Common Response Patterns}\tabularnewline
\toprule\noalign{}
Pattern & Freq \\
\midrule\noalign{}
\endfirsthead
\toprule\noalign{}
Pattern & Freq \\
\midrule\noalign{}
\endhead
11111 & 298 \\
11011 & 173 \\
10011 & 81 \\
10111 & 80 \\
11101 & 61 \\
11001 & 56 \\
10001 & 29 \\
10101 & 28 \\
\bottomrule\noalign{}
\end{longtable}
\end{frame}

\begin{frame}{Initial Setup}
\phantomsection\label{initial-setup}
\textbf{2PL model} with intercept/slope parameterization:

\[P_i(\theta) = \frac{1}{1 + \exp(-(d_i + a_i\theta))}\]

\textbf{Starting values:} all intercepts \(d_i = 0\), all slopes
\(a_i = 1\)

\textbf{Quadrature:} 10 nodes from \(-4\) to \(+4\) with weights from
the standard normal

These starting values treat all items as identical --- the \emph{data}
will differentiate them.
\end{frame}

\begin{frame}{E-Step 1: Likelihood at Each Node}
\phantomsection\label{e-step-1-likelihood-at-each-node}
For each response pattern \(l\) at each quadrature node \(X_k\):

\[L_l(X_k) = \prod_{i=1}^{n} P_i(X_k)^{u_{li}} Q_i(X_k)^{1-u_{li}}\]

``How probable is this response pattern if the examinee's ability were
exactly \(X_k\)?''

\begin{longtable}[]{@{}cc@{}}
\caption{Pattern 22 \{1,0,1,0,1\}: Likelihood at each
node}\tabularnewline
\toprule\noalign{}
X\_k & L\_l(X\_k) \\
\midrule\noalign{}
\endfirsthead
\toprule\noalign{}
X\_k & L\_l(X\_k) \\
\midrule\noalign{}
\endhead
-4.0000 & 0.00000561 \\
-3.1111 & 0.00007111 \\
-2.2222 & 0.00076087 \\
-1.3333 & 0.00568600 \\
-0.4444 & 0.02214009 \\
0.4444 & 0.03452866 \\
1.3333 & 0.02157010 \\
2.2222 & 0.00702097 \\
3.1111 & 0.00159618 \\
4.0000 & 0.00030636 \\
\bottomrule\noalign{}
\end{longtable}
\end{frame}

\begin{frame}{E-Step 2: Marginal Probability}
\phantomsection\label{e-step-2-marginal-probability}
Sum over nodes, weighted by quadrature weights:

\[P_l = \sum_{k=1}^{q} L_l(X_k) \cdot A(X_k)\]

This is the \textbf{marginal likelihood} of pattern \(l\) --- the
probability of seeing that pattern averaged across all ability levels.

\begin{longtable}[]{@{}lrr@{}}
\toprule\noalign{}
Pattern & Freq & Marginal Prob \\
\midrule\noalign{}
\endhead
11111 & 298 & 0.0955352 \\
11011 & 173 & 0.0360619 \\
10011 & 81 & 0.0224199 \\
10111 & 80 & 0.0360619 \\
11101 & 61 & 0.0360619 \\
11001 & 56 & 0.0224199 \\
10001 & 29 & 0.0224199 \\
10101 & 28 & 0.0224199 \\
\bottomrule\noalign{}
\end{longtable}
\end{frame}

\begin{frame}{E-Step 3: Posterior Probabilities}
\phantomsection\label{e-step-3-posterior-probabilities}
Apply Bayes' theorem at each node:

\[P(X_k \mid \mathbf{u}_l) = \frac{L_l(X_k) \cdot A(X_k)}{P_l}\]

\begin{longtable}[]{@{}ccccc@{}}
\caption{Pattern 22 \{1,0,1,0,1\}; \(P_l\) = 0.02241993}\tabularnewline
\toprule\noalign{}
\(X_k\) & \(L_l(X_k)\) & \(A(X_k)\) & \(L \times A\) & Posterior \\
\midrule\noalign{}
\endfirsthead
\toprule\noalign{}
\(X_k\) & \(L_l(X_k)\) & \(A(X_k)\) & \(L \times A\) & Posterior \\
\midrule\noalign{}
\endhead
-4.0000 & 0.00000561 & 0.000119 & 0.00000000 & 0.00000003 \\
-3.1111 & 0.00007111 & 0.002805 & 0.00000020 & 0.00000890 \\
-2.2222 & 0.00076087 & 0.030020 & 0.00002284 & 0.00101879 \\
-1.3333 & 0.00568600 & 0.145800 & 0.00082902 & 0.03697690 \\
-0.4444 & 0.02214009 & 0.321300 & 0.00711361 & 0.31728951 \\
0.4444 & 0.03452866 & 0.321300 & 0.01109406 & 0.49483018 \\
1.3333 & 0.02157010 & 0.145800 & 0.00314492 & 0.14027339 \\
2.2222 & 0.00702097 & 0.030020 & 0.00021077 & 0.00940098 \\
3.1111 & 0.00159618 & 0.002805 & 0.00000448 & 0.00019970 \\
4.0000 & 0.00030636 & 0.000119 & 0.00000004 & 0.00000163 \\
\bottomrule\noalign{}
\end{longtable}
\end{frame}

\begin{frame}{E-Step 3: Posterior (Visual)}
\phantomsection\label{e-step-3-posterior-visual}
\begin{center}\includegraphics{MMLE_EM_Slides_files/figure-beamer/e3-plot-1} \end{center}
\end{frame}

\begin{frame}{E-Step 4: The ``Artificial Data''}
\phantomsection\label{e-step-4-the-artificial-data}
Aggregate across all examinees (weighted by pattern frequency \(f_l\)):

\[\bar{n}_k = \sum_{l=1}^{s} f_l \cdot P(X_k \mid \mathbf{u}_l) \quad \text{(expected examinees at node } k\text{)}\]

\[\bar{r}_{ik} = \sum_{l=1}^{s} u_{li} \cdot f_l \cdot P(X_k \mid \mathbf{u}_l) \quad \text{(expected correct on item } i \text{ at node } k\text{)}\]

These play the same role as the observed frequencies in the JMLE
estimation equations.
\end{frame}

\begin{frame}{E-Step 4: Artificial Data Table}
\phantomsection\label{e-step-4-artificial-data-table}
\begin{longtable}[]{@{}
  >{\centering\arraybackslash}p{(\linewidth - 12\tabcolsep) * \real{0.0882}}
  >{\centering\arraybackslash}p{(\linewidth - 12\tabcolsep) * \real{0.1275}}
  >{\centering\arraybackslash}p{(\linewidth - 12\tabcolsep) * \real{0.1569}}
  >{\centering\arraybackslash}p{(\linewidth - 12\tabcolsep) * \real{0.1569}}
  >{\centering\arraybackslash}p{(\linewidth - 12\tabcolsep) * \real{0.1569}}
  >{\centering\arraybackslash}p{(\linewidth - 12\tabcolsep) * \real{0.1569}}
  >{\centering\arraybackslash}p{(\linewidth - 12\tabcolsep) * \real{0.1569}}@{}}
\caption{Artificial data from E-Step (Cycle 1)}\tabularnewline
\toprule\noalign{}
\begin{minipage}[b]{\linewidth}\centering
\(X_k\)
\end{minipage} & \begin{minipage}[b]{\linewidth}\centering
\(\bar{n}_k\)
\end{minipage} & \begin{minipage}[b]{\linewidth}\centering
\(\bar{r}_{1k}\)
\end{minipage} & \begin{minipage}[b]{\linewidth}\centering
\(\bar{r}_{2k}\)
\end{minipage} & \begin{minipage}[b]{\linewidth}\centering
\(\bar{r}_{3k}\)
\end{minipage} & \begin{minipage}[b]{\linewidth}\centering
\(\bar{r}_{4k}\)
\end{minipage} & \begin{minipage}[b]{\linewidth}\centering
\(\bar{r}_{5k}\)
\end{minipage} \\
\midrule\noalign{}
\endfirsthead
\toprule\noalign{}
\begin{minipage}[b]{\linewidth}\centering
\(X_k\)
\end{minipage} & \begin{minipage}[b]{\linewidth}\centering
\(\bar{n}_k\)
\end{minipage} & \begin{minipage}[b]{\linewidth}\centering
\(\bar{r}_{1k}\)
\end{minipage} & \begin{minipage}[b]{\linewidth}\centering
\(\bar{r}_{2k}\)
\end{minipage} & \begin{minipage}[b]{\linewidth}\centering
\(\bar{r}_{3k}\)
\end{minipage} & \begin{minipage}[b]{\linewidth}\centering
\(\bar{r}_{4k}\)
\end{minipage} & \begin{minipage}[b]{\linewidth}\centering
\(\bar{r}_{5k}\)
\end{minipage} \\
\midrule\noalign{}
\endhead
-4.0000 & 0.0 & 0.0 & 0.0 & 0.0 & 0.0 & 0.0 \\
-3.1111 & 0.1 & 0.0 & 0.0 & 0.0 & 0.0 & 0.0 \\
-2.2222 & 2.7 & 1.4 & 0.4 & 0.2 & 0.5 & 1.0 \\
-1.3333 & 31.1 & 22.1 & 9.6 & 4.9 & 11.7 & 18.0 \\
-0.4444 & 182.0 & 155.1 & 91.1 & 56.3 & 105.5 & 137.5 \\
0.4444 & 402.3 & 374.1 & 278.7 & 206.0 & 304.1 & 351.3 \\
1.3333 & 295.8 & 286.4 & 249.0 & 211.1 & 259.6 & 278.2 \\
2.2222 & 77.6 & 76.5 & 72.0 & 66.5 & 73.3 & 75.5 \\
3.1111 & 8.1 & 8.0 & 7.8 & 7.6 & 7.9 & 8.0 \\
4.0000 & 0.4 & 0.4 & 0.4 & 0.3 & 0.4 & 0.4 \\
\bottomrule\noalign{}
\end{longtable}
\end{frame}

\begin{frame}[fragile]{Why Do the \(\bar{r}_{ik}\) Columns Differ?}
\phantomsection\label{why-do-the-barr_ik-columns-differ}
We started with \textbf{identical} parameters for all 5 items, yet the
\(\bar{r}_{ik}\) values differ.

\begin{itemize}
\tightlist
\item
  \(\bar{n}_k\) depends only on the posterior --- \textbf{same for all
  items}
\item
  \(\bar{r}_{ik}\) multiplies by \(u_{li}\) (the actual response) ---
  only correct responses contribute
\end{itemize}

\textbf{The raw data break the symmetry:}

\begin{verbatim}
##   Item 1: 0.924 proportion correct
##    Item 2: 0.709 proportion correct
##    Item 3: 0.553 proportion correct
##    Item 4: 0.763 proportion correct
##    Item 5: 0.87 proportion correct
\end{verbatim}

Items with higher proportion correct \(\rightarrow\) larger
\(\bar{r}_{ik}\) values.
\end{frame}

\begin{frame}{The M-Step: Estimating Item Parameters}
\phantomsection\label{the-m-step-estimating-item-parameters}
Treat the artificial data as known and solve (one item at a time):

\[\sum_{k=1}^{q} \bar{n}_k \left[\frac{\bar{r}_{ik}}{\bar{n}_k} - P_i(X_k)\right] = 0 \quad \text{(for } d_i\text{)}\]

\[\sum_{k=1}^{q} \bar{n}_k \left[\frac{\bar{r}_{ik}}{\bar{n}_k} - P_i(X_k)\right] X_k = 0 \quad \text{(for } a_i\text{)}\]

\textbf{Logic:} Find the \(d_i\) and \(a_i\) that make the residuals
(observed \(-\) expected) sum to zero. The \(X_k\) multiplier in the
second equation captures the \emph{trend} across ability (slope).

Solve iteratively using \textbf{Newton-Raphson} (same as in JMLE).
\end{frame}

\begin{frame}{M-Step: Results After Cycle 1}
\phantomsection\label{m-step-results-after-cycle-1}
\begin{longtable}[]{@{}cccc@{}}
\caption{Item Parameter Estimates After 1 EM Cycle}\tabularnewline
\toprule\noalign{}
Item & Intercept (\(d_i\)) & Slope (\(a_i\)) & Difficulty (\(b_i\)) \\
\midrule\noalign{}
\endfirsthead
\toprule\noalign{}
Item & Intercept (\(d_i\)) & Slope (\(a_i\)) & Difficulty (\(b_i\)) \\
\midrule\noalign{}
\endhead
1 & 2.1659 & 0.9438 & -2.2949 \\
2 & 0.4102 & 0.9410 & -0.4360 \\
3 & -0.3787 & 0.9697 & 0.3905 \\
4 & 0.7263 & 0.9322 & -0.7791 \\
5 & 1.5321 & 0.9198 & -1.6658 \\
\bottomrule\noalign{}
\end{longtable}

Already, the items have differentiated! Items with higher proportion
correct get larger (more negative) difficulty estimates.
\end{frame}

\begin{frame}{The Full EM Loop}
\phantomsection\label{the-full-em-loop}
Now iterate: use these new estimates as input to the next E-step,
repeat.

\begin{longtable}[]{@{}cccc@{}}
\caption{Final Estimates After 44 EM Cycles}\tabularnewline
\toprule\noalign{}
Item & Intercept (\(d_i\)) & Slope (\(a_i\)) & Difficulty (\(b_i\)) \\
\midrule\noalign{}
\endfirsthead
\toprule\noalign{}
Item & Intercept (\(d_i\)) & Slope (\(a_i\)) & Difficulty (\(b_i\)) \\
\midrule\noalign{}
\endhead
1 & 2.7732 & 0.8256 & -3.3588 \\
2 & 0.9903 & 0.7231 & -1.3695 \\
3 & 0.2491 & 0.8901 & -0.2798 \\
4 & 1.2848 & 0.6887 & -1.8657 \\
5 & 2.0535 & 0.6573 & -3.1239 \\
\bottomrule\noalign{}
\end{longtable}
\end{frame}

\begin{frame}{Watching EM Converge}
\phantomsection\label{watching-em-converge}
\begin{center}\includegraphics{MMLE_EM_Slides_files/figure-beamer/convergence-d-1} \end{center}
\end{frame}

\begin{frame}{Watching EM Converge (cont.)}
\phantomsection\label{watching-em-converge-cont.}
\begin{center}\includegraphics{MMLE_EM_Slides_files/figure-beamer/convergence-a-1} \end{center}

EM converges \textbf{slowly} --- parameters change rapidly early, then
creep toward final values.
\end{frame}

\begin{frame}[fragile]{Comparison with \texttt{mirt}}
\phantomsection\label{comparison-with-mirt}
\begin{longtable}[]{@{}ccccc@{}}
\caption{Our MMLE/EM vs.~mirt Package}\tabularnewline
\toprule\noalign{}
Item & Our \(a\) & mirt \(a\) & Our \(b\) & mirt \(b\) \\
\midrule\noalign{}
\endfirsthead
\toprule\noalign{}
Item & Our \(a\) & mirt \(a\) & Our \(b\) & mirt \(b\) \\
\midrule\noalign{}
\endhead
1 & 0.826 & 0.825 & -3.359 & -3.361 \\
2 & 0.723 & 0.723 & -1.369 & -1.370 \\
3 & 0.890 & 0.890 & -0.280 & -0.280 \\
4 & 0.689 & 0.689 & -1.866 & -1.866 \\
5 & 0.657 & 0.658 & -3.124 & -3.123 \\
\bottomrule\noalign{}
\end{longtable}

Small differences reflect implementation details: \texttt{mirt} uses 61
quadrature points (vs.~our 10) and Bock-Aitkin acceleration.
\end{frame}

\begin{frame}[fragile]{EM Algorithm Summary}
\phantomsection\label{em-algorithm-summary}
\begin{verbatim}
1. INITIALIZE: Starting values for all item parameters

2. E-STEP:
   → Compute likelihood at each node for each response pattern
   → Apply Bayes' theorem to get posterior at each node
   → Aggregate into "artificial data" (n-bar, r-bar)

3. M-STEP:
   → For each item separately:
     Treat artificial data as known, solve for d and a
     using Newton-Raphson

4. CONVERGED? If yes → stop. If no → go to step 2.
\end{verbatim}
\end{frame}

\begin{frame}{Key Takeaways}
\phantomsection\label{key-takeaways}
\begin{itemize}
\item
  \textbf{JMLE} estimates every person's \(\theta\) along with item
  parameters \(\rightarrow\) inconsistent estimates (Neyman-Scott
  problem)
\item
  \textbf{MMLE} integrates out \(\theta\) by assuming a population
  distribution \(\rightarrow\) consistent item parameter estimates
\item
  \textbf{EM} handles the circular dependency: the E-step infers the
  ability distribution, the M-step re-estimates item parameters, repeat
\item
  MMLE produces \textbf{item parameters only} --- person \(\theta\)
  estimates require a separate step (MLE, MAP, or EAP)
\end{itemize}
\end{frame}

\end{document}
